\documentclass[11pt]{article}
\usepackage[utf8]{inputenc}
\usepackage[T1]{fontenc}
\usepackage{amsmath,amssymb}
\usepackage{graphicx}
\usepackage{booktabs}
\usepackage[margin=1in]{geometry}
\usepackage{hyperref}

\title{Verificacion de Gibbs, Parametro Beta, y Experimentos a Escala\\
\large (Notebook: phase3\_findings)}
\author{}
\date{}

\begin{document}
\maketitle

\begin{abstract}
Este documento cubre tres temas:
(1) encontramos y corregimos un bug en el muestreador de Gibbs,
(2) exploramos el parametro $\beta$ que controla cuanto le importa
al algoritmo ``aprender'' vs ``aclarar'', y
(3) corrimos experimentos preliminares con 20 personas.
Los parametros $\alpha$ y las comparaciones con la tesis de Nico
estan en el documento 03.
\end{abstract}

%% ====================================================================
\section{El Bug del Gibbs y Como lo Arreglamos}

\subsection{Que es el Gibbs?}

Cuando tenemos muchas personas ($n > 20$), no podemos calcular las
probabilidades exactas de infeccion despues de un test.
Usamos \textbf{muestreo de Gibbs}: un metodo que genera muchas
combinaciones posibles de ``quien esta infectado'' y promedia sobre ellas.

\subsection{El problema}

Encontramos que Gibbs daba respuestas incorrectas cuando los pools
se traslapaban (compartian personas).

\textbf{Ejemplo:} si pool 1 = $\{A, B\}$ y pool 2 = $\{B, C\}$,
los resultados de ambos tests restringen juntos quien puede estar infectado.
Gibbs no exploraba bien todas las combinaciones validas --- se quedaba
``atrapado'' en una region del espacio.

\textbf{Antes del fix:} 85.8\% de tests pasaban.
\textbf{Despues del fix:} 100\% de tests pasaban (64/64).

\subsection{La solucion}

Dos cambios:
\begin{enumerate}
  \item Si quedan $\leq 7$ personas activas, calcular exacto (no Gibbs).
  \item Agregar ``movimientos de Metropolis-Hastings'': proponer intercambios
        globales de infectado/sano entre pares de personas.
        Esto permite saltar entre regiones desconectadas.
\end{enumerate}

\subsection{Verificacion de convergencia}

Probamos con distintas cantidades de iteraciones:
50, 100, 200, 500, 1000, 2000, 5000.

\begin{itemize}
  \item Con $\geq 500$ iteraciones, el error maximo cae por debajo de 0.03.
  \item Con $\geq 2000$, el error es $< 0.01$.
  \item Conclusion: 1000 iteraciones es un buen default.
\end{itemize}

%% ====================================================================
\section{El Parametro Beta ($\beta$): Aclarar vs Aprender}

\subsection{Que es $\beta$?}

Cuando el algoritmo greedy elige que grupo testear, normalmente solo
le importa una cosa: \emph{aclarar personas} (que el resultado sea $r=0$).

Su puntaje para cada pool es:
$$\text{score}_{\text{normal}} = P(r=0) \times \sum u_i$$

Es decir: ``la probabilidad de que todos en el pool esten sanos, multiplicada
por el valor total de aclarar a esas personas''.

Pero hay otra cosa util: \emph{aprender quien esta infectado}.
Aunque $r > 0$ (el pool sale positivo), el numero exacto $r$ nos da
informacion valiosa sobre las personas del pool.

$\beta$ agrega un bonus por informacion:
$$\text{score}_{\beta} = P(r=0) \times \sum u_i \;+\; \beta \times \text{ganancia\_de\_info}$$

\begin{itemize}
  \item $\beta = 0$: greedy normal. Solo quiere aclarar.
  \item $\beta > 0$: tambien valora aprender. Puede elegir pools mas grandes
        o mas ``informativos'' aunque tengan menor chance de salir negativos.
\end{itemize}

\textbf{Analogia:} es como decidir entre una pregunta de examen facil
(que sabes que vas a contestar bien) vs una dificil (que te ensena mas
para las siguientes preguntas).

\subsection{Que mide la ``ganancia de informacion''?}

Usamos la metrica \texttt{confirmed}: cuenta cuantas personas
quedarian con probabilidad $> 0.95$ o $< 0.05$ despues del test.
Es decir, cuantas personas ``resolvemos'' en terminos de certeza,
aunque no las ``aclaremos'' formalmente.

\subsection{Experimento 1: Prevalencia alta ($p \approx 0.8$)}

Escenario VIP con $n=20$, $B=6$, $G=10$. Los VIPs tienen $p=0.8$.

\textbf{Resultado:} $\beta$ tiene efecto limitado.

Por que? Porque con $p=0.8$, la probabilidad de que un pool salga
negativo es casi cero ($0.2^k$ para un pool de tamano $k$).
El score base es muy bajo, asi que $\beta$ tiene efecto limitado
a valores chicos.
Con $\beta$ grande (2+), la EU puede subir moderadamente
y el algoritmo cambia a pools mas grandes, pero la relacion no es monotona.

\subsection{Experimento 2: Prevalencia moderada ($p \approx 0.35$)}

Escenario VIP similar con $n=20$, $B=6$, $G=5$, VIPs con $p=0.35$.

\textbf{Resultado:} $\beta$ tiene efecto moderado.
\begin{itemize}
  \item $\beta = 0$: pools chicos, EU base.
  \item $\beta = 2$: pools mas grandes, mejor EU entre los valores probados.
  \item $\beta = 1$: no necesariamente mejor que $\beta = 0$ --- la relacion no es monotona.
\end{itemize}

\subsection{Resumen de $\beta$}

\begin{table}[htbp]
  \centering
  \begin{tabular}{lcc}
    \toprule
    Prevalencia & Efecto de $\beta$ & Recomendacion \\
    \midrule
    Baja ($p < 0.1$) & No probado & Probablemente innecesario (pools ya funcionan bien) \\
    Moderada ($p \approx 0.2$--$0.4$) & Moderado & Probar $\beta \in [1, 3]$ \\
    Alta ($p > 0.5$) & Limitado & $\beta$ no ayuda mucho \\
    \bottomrule
  \end{tabular}
  \caption{Cuando usar $\beta$.}
\end{table}

%% ====================================================================
\section{Experimentos Preliminares a Gran Escala}

Corrimos experimentos con $n=20$, $B=2$, $G=10$ (configuracion D).
Estos son resultados \textbf{preliminares} con 1--2 instancias por configuracion
(se necesitan $\sim$50 para conclusiones estadisticas).

\subsection{Resultados}

\begin{itemize}
  \item El greedy (con Mosek para seleccion de pools) captura $\sim$35\% del
        gap entre $U_{\text{single}}$ y $U_{\max}$.
  \item El gap es grande: con solo $B=2$ tests para 20 personas,
        hay mucho margen de mejora.
  \item Con escenario VIP ($B=6$): el gap se reduce.
  \item Con utilidades sesgadas (algunos valen mucho mas): el gap tambien cambia.
\end{itemize}

\textbf{Limitacion:} estos son quick-runs. Los numeros exactos pueden cambiar
con mas instancias.

\end{document}
